\cgeentry{cge_sensitivity}{Sensitivity of a scenario's results to the elasticities}

\cgesec{Description}

\cmd{cge_sensitivity} shows how much the results of a scenario depend on the
elasticities. For each of a number of draws, it multiplies the elasticities
by random factors around one, calibrates the model again from the SAM,
restores the closure, runs a program that applies the scenario's shocks, and
solves. It then summarises every report measure of \manref{cge_report} across
the draws that converged: the result with the central elasticities, the
mean, the standard deviation and the 5th and 95th percentiles. The model in
memory, with its simulations, is restored exactly at the end.

\cgesec{Quick start}

200 draws of every value-added, Armington and CET elasticity within 50 per
cent of its central value, for the shocks applied by program \cmd{tariff_cut}
\begin{stcode}
cge_sensitivity, scenario(tariff_cut) draws(200)
\end{stcode}
As above, with factors drawn within 25 per cent from a triangular distribution,
and every report table displayed
\begin{stcode}
cge_sensitivity, scenario(tariff_cut) draws(200) spread(0.25) distribution(triangular) tables(all)
\end{stcode}
Vary only the Armington and CET elasticities, with one factor for each type
applied to every commodity, and load every draw as data
\begin{stcode}
cge_sensitivity, scenario(tariff_cut) draws(200) vary(q t) common clear
\end{stcode}
Save every draw to \cmd{sens.dta} without displaying the summary
\begin{stcode}
cge_sensitivity, scenario(tariff_cut) draws(200) saving(sens, replace) nodisplay
\end{stcode}

\cgesec{Syntax}

\begin{stsyntax}
\cmdb{cge_sensitivity}\texttt{,} \optarg{scen}{ario}{program} \optarg{draw}{s}{\#}
[\argu{options}]
\end{stsyntax}

\begin{opttable}
\optgroup{Main}
\optarg{scen}{ario}{program} & program that applies the scenario's shocks; required\\
\optarg{draw}{s}{\#} & number of draws; required\\
\optarg{}{seed}{\#} & random-number seed; default \texttt{12345}\\
\optgroup{Draws}
\optarg{spr}{ead}{\#} & factors drawn between $1-\#$ and $1+\#$; default \texttt{0.5}\\
\optarg{dist}{ribution}{name} & \cmd{uniform} or \cmd{triangular}; default \cmd{uniform}\\
\opt{}{common} & one factor for each type of elasticity, applied to every sector\\
\optarg{var}{y}{list} & elasticities to vary: \cmd{va}, \cmd{q}, \cmd{t}, \cmd{a}; default \cmd{va q t}\\
\optgroup{Reporting}
\optarg{tab}{les}{list} & report tables to display; default \cmd{macro welfare}\\
\opt{nod}{isplay} & suppress the summary\\
\opt{}{clear} & replace the data in memory with every draw\\
\optarg{sav}{ing}{filename\textup{[}\texttt{, replace}\textup{]}} & save every draw to a Stata dataset\\
\optarg{}{keep}{name} & keep every draw as a simulation\\
\end{opttable}

\cgesec{Options}

\begin{optdesc}
\item[\optarg{scen}{ario}{program}] names a program, without arguments, that
applies the scenario's shocks with \manref{cge_shock} to the freshly
calibrated model. It is run quietly for every draw. If it fails, the model is
restored and the command stops.

\item[\optarg{draw}{s}{\#}] sets the number of draws, at least 1. Draw 0, with
the central elasticities, is always added, so the model is calibrated and
solved \# + 1 times.

\item[\optarg{}{seed}{\#}] sets the random-number seed; the default is 12345.
It sets the seed of Stata's random-number generator. The same seed gives the
same draws, whichever elasticities are varied.

\item[\optarg{spr}{ead}{\#}] sets the range of the factors, from $1 - \#$ to
$1 + \#$; the default is 0.5. \# must be at least 0 and less than 1, so
that every drawn elasticity keeps its sign.

\item[\optarg{dist}{ribution}{name}] draws the factors from a
\cmd{uniform} distribution on that range, the default, or from a symmetric
\cmd{triangular} distribution on the same range, which puts more weight near
the central values.

\item[\opt{}{common}] draws one factor for each type of elasticity and each
draw, applied to every sector, so that the elasticities of a type move
together. By default, each sector's elasticity is drawn separately.

\item[\optarg{var}{y}{list}] lists the elasticities to vary: \cmd{va}, the
elasticity of substitution in value added of each activity; \cmd{q}, the
Armington elasticity of each commodity; \cmd{t}, the CET elasticity of each
commodity, which also applies to exports drawn from composite supply; and
\cmd{a}, the elasticity of the top nest of each activity. The default is
\cmd{va q t}. A Leontief top nest, with elasticity 0, stays Leontief. Other
elasticities and calibration settings keep their values.

\item[\optarg{tab}{les}{list}] lists the tables of \manref{cge_report} whose
measures are displayed in the summary, or \cmd{all}; the default is
\cmd{macro welfare}. The dataset of \opt{}{clear} and \opt{sav}{ing()} holds
every table.

\item[\opt{nod}{isplay}] suppresses the summary. The progress of the draws is
still displayed.

\item[\opt{}{clear}] replaces the data in memory with the results described
under Remarks and examples below.

\item[\optarg{sav}{ing}{filename\textup{[}\texttt{, replace}\textup{]}}]
saves the same dataset as a Stata dataset, leaving the data in memory
unchanged unless \opt{}{clear} is also given; \cmd{replace} allows an existing
file to be replaced.

\item[\optarg{}{keep}{name}] (from version 0.8.3) keeps the solution of every
draw that converged as a simulation of the model: \argu{name}\cmd{0}, with the
central elasticities, \argu{name}\cmd{1}, \argu{name}\cmd{2}, and so on, each
with the parameters of its draw. \manref{cge_microsim} with
\cmd{bands(}\argu{name}\cmd{)} then gives the range of poverty and inequality
across the draws, and any command that takes a simulation, such as
\manref{cge_results}, can use them. Simulations of the same names are
replaced.
\end{optdesc}

\cgesec{Remarks and examples}

Elasticities are rarely estimated for the country and period of a SAM, and a
scenario's results can depend on them. \cmd{cge_sensitivity} solves the
scenario again for many sets of elasticities around those used. Each draw
calibrates the model from the SAM with the drawn elasticities and the other
settings of the last \manref{cge_calibrate}, so that every draw replicates
the SAM; it then restores the closure, including swaps, and any constraints,
runs the scenario program and solves. A model calibrated by a version of
\pkg that did not record its calibration settings must be calibrated again
first (error 459). The solution of each draw is
summarised by the report measures of \manref{cge_report}: percentage changes
from the base, changes in points for rates, and per cent of base consumption
for the welfare measures.

The draws do not estimate uncertainty about the elasticities. They show how
the results change over a range chosen by the user, so the range, the
distribution and whether sectors move together should be stated with the
results (\chapref{chap:sensitivity}). Standard errors of published estimates,
such as those of \mbox{\manref{cge_elasticities}}, can guide the choice of range.

\cgesubsec{Example 1: the elasticities of every sector drawn separately}

The synthetic SAM is calibrated as in \manref{cge_solve}, with
\cmd{sigmava(0.8) sigmaq(2) sigmat(2)}. The scenario removes every import
tariff; its program applies the shock and nothing else:
\stlog{sens_program}
Twenty draws, each elasticity of each sector drawn uniformly within 50 per
cent of its central value:
\stlog{sens_draws}
The column \cmd{central} is the result of the scenario with the central
elasticities, the result of \mbox{\manref{cge_solve}}; the other columns summarise
the 20 draws. Real consumption and household welfare vary little across the
draws, by a standard deviation of about 0.02 points around 2.7 per cent; real
exports and imports, and the exchange rate, vary much more. With 20 draws,
\cmd{p5} is the smallest result and \cmd{p95} the second largest; with more
draws the percentiles are more precise. The results come from a synthetic SAM and
illustrative elasticities and are not findings about any economy.

\cgesubsec{Example 2: common factors, and the draws as data}

With \opt{}{common}, the Armington and CET elasticities of every commodity
are multiplied by the same factor in each draw, here drawn from a triangular
distribution; only the welfare table is displayed, and every draw is loaded:
\stlog{sens_common}
The dataset holds the report measures of each draw that converged, draw 0
included, and the elasticities of every draw (table \cmd{elasticity}). With
\opt{}{common}, the Armington elasticities are the same for every commodity
in each draw:
\stlog{sens_data}
The mean and standard deviation of the equivalent variation over draws 1 to
20 are those of the summary.

\cgesubsec{The dataset}

The dataset of \opt{}{clear} and \opt{sav}{ing()} has the variables
\begin{cgetable}{>{\ttfamily}l p{0.80\linewidth}}
\toprule
\normalfont Variable & Content\\
\midrule\endhead
\bottomrule\endlastfoot
draw & the draw; 0 for the central elasticities\\
table & the report table, or \cmd{elasticity}\\
measure & the report measure, or \cmd{sigmava}, \cmd{sigmaq}, \cmd{sigmat} or \cmd{sigmaa}\\
element & the set element\\
base & the base value of the measure\\
value & its value in the draw, or the value of the elasticity\\
pct & the per cent change from the base (rates: points; welfare: per cent of base consumption)\\
\end{cgetable}
Report measures are included for the draws that converged; elasticities for
every draw, with \cmd{base} and \cmd{pct} missing. A draw that does not
calibrate or does not solve is counted as not converged and left out of the
summary. If the scenario does not solve with the central elasticities, the
command stops with error 430 after restoring the model.

\cgesubsec{The model in memory}

The model is saved before the first draw and loaded again after the last, so
the elasticities, closure, shocks and simulations in memory are those before
the command. The simulation that each draw solves is not kept. The time taken
is about \argu{\#} + 1 times that of a calibration and a solve.

\cgesec{Stored results}

\cmd{cge_sensitivity} stores the following in \cmd{r()}:
\begin{storedres}
\resgroup{Scalars}
\res{r(N_draws)} & number of draws\\
\res{r(N_converged)} & number of draws that converged\\
\end{storedres}

\cgesec{Methods and formulas}

Let $\sigma^0_k$ be the central value of elasticity $k$ of a varied type,
and $s$ the value of \cmd{spread()}. In draw $d = 1, \dots, D$,
\begin{equation*}
\sigma^d_k = \sigma^0_k\,(1 + s\,u_{kd}), \qquad
u_{kd} =
\begin{cases}
2\,U_{kd} - 1 & \text{uniform},\\
U_{kd} + U'_{kd} - 1 & \text{triangular},
\end{cases}
\end{equation*}
where $U_{kd}$ and $U'_{kd}$ are independent uniform draws on $(0, 1)$, so
that $u_{kd}$ is uniform, or symmetric triangular with density $1 - |u|$, on
$(-1, 1)$. With \opt{}{common}, $u_{kd} = u_d$ for every element $k$ of a
type. The draws are made after the seed is set, for the four types in a
fixed order (value added, Armington, CET, top nest), whether or not a type is
varied, so that a seed gives the same factors for a type whichever types are
varied. $\sigma^0_k = 0$ gives $\sigma^d_k = 0$.

For a report measure with results $y_d$ in the set $S$ of draws $d \ge 1$
that converged, $n = |S|$, the summary gives the central result $y_0$,
\begin{equation*}
\bar y = \frac{1}{n}\sum_{d \in S} y_d, \qquad
s_y = \Bigl[\frac{1}{n-1}\sum_{d \in S} (y_d - \bar y)^2\Bigr]^{1/2},
\end{equation*}
and the order statistics $y_{(\lceil 0.05\,n\rceil)}$ and
$y_{(\lceil 0.95\,n\rceil)}$ as the 5th and 95th percentiles, with
$y_{(1)} \le \dots \le y_{(n)}$. \Chapref{chap:sensitivity} discusses the method
and its interpretation.

\cgesec{Also see}

\begin{alsosee}
\seeentry{cge_calibrate}{Calibrate the model to the SAM}
\seeentry{cge_report}{Standard report tables for one or more simulations}
\seeentry{cge_elasticities}{Published elasticities, with their sources}
\seeentry{cge_shock}{Change a parameter or a fixed variable}
\seechapter{chap:sensitivity}{Sensitivity analysis}
\end{alsosee}
